\documentclass[12pt, a4paper]{article}

% ─── Codificación y lenguaje ───────────────────────────────────────────────────
\usepackage[utf8]{inputenc}
\usepackage[T1]{fontenc}
\usepackage[spanish, es-noshorthands]{babel}

% ─── Fuentes ──────────────────────────────────────────────────────────────────
\usepackage{lmodern}
\usepackage{microtype}

% ─── Matemáticas ──────────────────────────────────────────────────────────────
\usepackage{amsmath, amssymb, amsthm, mathtools}

% ─── Geometría y diseño ───────────────────────────────────────────────────────
\usepackage[
    top=2.5cm, bottom=2.5cm,
    left=3cm,  right=3cm
]{geometry}
\usepackage{parskip}
\usepackage{setspace}
\onehalfspacing

% ─── Color y gráficos ─────────────────────────────────────────────────────────
\usepackage[dvipsnames, table]{xcolor}
\definecolor{azulprincipal}{RGB}{25, 60, 115}
\definecolor{azulclaro}{RGB}{210, 225, 245}
\definecolor{grisclaro}{RGB}{245, 247, 250}
\definecolor{dorado}{RGB}{180, 140, 40}

% ─── Cajas y marcos ───────────────────────────────────────────────────────────
\usepackage{tcolorbox}
\tcbuselibrary{skins, breakable, theorems}

% Caja para teoremas
\tcbset{
    teostyle/.style={
        enhanced, breakable,
        colback=azulclaro!30,
        colframe=azulprincipal,
        fonttitle=\bfseries\color{azulprincipal},
        attach boxed title to top left={yshift=-2mm, xshift=6mm},
        boxed title style={
            colback=azulprincipal,
            colframe=azulprincipal,
            sharp corners
        },
        sharp corners=south,
        left=6pt, right=6pt, top=8pt, bottom=6pt,
    },
    resultstyle/.style={
        enhanced,
        colback=grisclaro,
        colframe=dorado,
        fonttitle=\bfseries\color{dorado},
        sharp corners,
        left=8pt, right=8pt, top=6pt, bottom=6pt,
        boxrule=1.2pt,
    },
    defstyle/.style={
        enhanced,
        colback=white,
        colframe=azulprincipal!50,
        fonttitle=\bfseries\color{azulprincipal},
        sharp corners,
        left=8pt, right=8pt, top=6pt, bottom=6pt,
        borderline west={3pt}{0pt}{azulprincipal},
        boxrule=0pt,
        frame hidden,
    }
}

% ─── Encabezados y pies de página ─────────────────────────────────────────────
\usepackage{fancyhdr}
\pagestyle{fancy}
\fancyhf{}
\fancyhead[L]{\small\color{azulprincipal}\textit{EDO con Transformadas de Laplace}}
\fancyhead[R]{\small\color{azulprincipal}\thepage}
\renewcommand{\headrulewidth}{0.6pt}
\renewcommand{\headrule}{\hbox to\headwidth{\color{azulprincipal}\leaders\hrule height \headrulewidth\hfill}}

% ─── Títulos de sección ────────────────────────────────────────────────────────
\usepackage{titlesec}
\titleformat{\section}
    {\large\bfseries\color{azulprincipal}}
    {\thesection.}{0.6em}{}[{\color{azulprincipal}\titlerule[0.6pt]}]
\titleformat{\subsection}
    {\normalsize\bfseries\color{azulprincipal!80}}
    {\thesubsection.}{0.5em}{}

% ─── Tabla de contenidos ──────────────────────────────────────────────────────
\usepackage{tocloft}
\renewcommand{\cftsecfont}{\bfseries\color{azulprincipal}}
\renewcommand{\cftsecpagefont}{\bfseries\color{azulprincipal}}
\renewcommand{\cftsubsecfont}{\color{azulprincipal!80}}
\renewcommand{\cftsubsecpagefont}{\color{azulprincipal!80}}

% ─── Hipervínculos ────────────────────────────────────────────────────────────
\usepackage[
    colorlinks=true,
    linkcolor=azulprincipal,
    urlcolor=azulprincipal
]{hyperref}

% ─── Comandos personalizados ──────────────────────────────────────────────────
\newcommand{\Lap}{\mathcal{L}}
\newcommand{\U}{\mathcal{U}}
\newcommand{\R}{\mathbb{R}}

% Entornos de teorema con tcolorbox
\newenvironment{teorema}[1]{%
    \begin{tcolorbox}[teostyle, title={Teorema: #1}]%
}{%
    \end{tcolorbox}%
}
\newenvironment{resultado}{%
    \begin{tcolorbox}[resultstyle]%
}{%
    \end{tcolorbox}%
}
\newenvironment{definicion}[1]{%
    \begin{tcolorbox}[defstyle, title={#1}]%
}{%
    \end{tcolorbox}%
}

% ══════════════════════════════════════════════════════════════════════════════
\begin{document}
% ══════════════════════════════════════════════════════════════════════════════

% ══════════════════════════════════════════════════════════════════════════════
\section{Planteamiento del problema}
% ══════════════════════════════════════════════════════════════════════════════

Se desea encontrar la función $y(t)$ que satisface el siguiente problema de valor inicial:

\begin{resultado}
\[
    y'' + y = f(t), \qquad y(0) = 1, \quad y'(0) = 0
\]
donde la función de forzamiento está definida por:
\[
    f(t) =
    \begin{cases}
        1, & 0 \leq t < \dfrac{\pi}{2} \\[8pt]
        \sin(t), & t \geq \dfrac{\pi}{2}
    \end{cases}
\]
\end{resultado}


% ══════════════════════════════════════════════════════════════════════════════
\section{Desarrollo de la solución}
% ══════════════════════════════════════════════════════════════════════════════

\subsection{Paso 1: Reescritura de \texorpdfstring{$f(t)$}{f(t)} mediante escalones}

El primer paso consiste en expresar la función $f(t)$ en términos de la función escalón unitario..

Dado que $f(t)$ cambia de valor en $t = \pi/2$, se escribe:
\[
    f(t) = 1 + \bigl(\sin(t) - 1\bigr)\,\U\!\left(t - \frac{\pi}{2}\right)
\]

\subsection{Paso 2: Aplicación de la transformada de Laplace}

Se aplica $\Lap\{\cdot\}$ a ambos lados de la ecuación diferencial:
\[
    \Lap\{y''\} + \Lap\{y\} = \Lap\{f(t)\}
\]

Usando las propiedades estándar y las condiciones iniciales $y(0) = 1$, $y'(0) = 0$:
\begin{align*}
    \Lap\{y\}   &= Y(s) \\
    \Lap\{y''\} &= s^2\,Y(s) - s\,y(0) - y'(0) = s^2\,Y(s) - s
\end{align*}

Para la transformada del forzamiento, se aplica el segundo teorema de traslación. Dado que el argumento del escalón es $t - \pi/2$, se evalúa $f$ desplazada:
\[
    \sin\!\left(\left(t - \frac{\pi}{2}\right) + \frac{\pi}{2}\right) - 1
    = \sin(t) - 1
    \quad \Longrightarrow \quad
    \text{el factor es }\sin(t) - 1\text{ en términos de }t
\]

Calculando por separado usando identidades trigonométricas ($\sin(t) = \cos(t - \pi/2)$):
\begin{align*}
    \Lap\{f(t)\}
    &= \Lap\{1\} + \Lap\!\left\{\bigl(\sin(t)-1\bigr)\,\U\!\left(t-\tfrac{\pi}{2}\right)\right\} \\[4pt]
    &= \frac{1}{s}
      + e^{-\frac{\pi}{2}s}\,\Lap\!\left\{\sin\!\left(t+\tfrac{\pi}{2}\right) - 1\right\} \\[4pt]
    &= \frac{1}{s}
      + e^{-\frac{\pi}{2}s}\,\Lap\{\cos(t) - 1\} \\[4pt]
    &= \frac{1}{s}
      + e^{-\frac{\pi}{2}s}\left(\frac{s}{s^2+1} - \frac{1}{s}\right)
\end{align*}

\subsection{Paso 3: Despeje de \texorpdfstring{$Y(s)$}{Y(s)}}

Sustituyendo en la ecuación transformada:
\[
    s^2\,Y - s + Y = \frac{1}{s} + \left(\frac{s}{s^2+1} - \frac{1}{s}\right)e^{-\frac{\pi}{2}s}
\]
\[
    Y\,(s^2 + 1) = s + \frac{1}{s} + \left(\frac{s}{s^2+1} - \frac{1}{s}\right)e^{-\frac{\pi}{2}s}
\]

Dividiendo entre $(s^2 + 1)$:
\[
    Y(s) = \frac{s}{s^2+1} + \frac{1}{s(s^2+1)}
           + \left(\frac{s}{(s^2+1)^2} - \frac{1}{s(s^2+1)}\right)e^{-\frac{\pi}{2}s}
\]

La solución en el dominio temporal se obtiene aplicando la transformada inversa a cada término:
\[
    y(t) = \Lap^{-1}\!\left\{\frac{s}{s^2+1}\right\}
         + \Lap^{-1}\!\left\{\frac{1}{s(s^2+1)}\right\}
         + \Lap^{-1}\!\left\{\frac{s\,e^{-\frac{\pi}{2}s}}{(s^2+1)^2}\right\}
         - \Lap^{-1}\!\left\{\frac{e^{-\frac{\pi}{2}s}}{s(s^2+1)}\right\}
\]

% ══════════════════════════════════════════════════════════════════════════════
\section{Cálculo de las transformadas inversas}
% ══════════════════════════════════════════════════════════════════════════════

\subsection{Término 1: \texorpdfstring{$\Lap^{-1}\!\left\{\dfrac{s}{s^2+1}\right\}$}{}}

Se aplica directamente la fórmula tabular de la transformada del coseno:

\begin{resultado}
\[
    \Lap^{-1}\!\left\{\frac{s}{s^2+1}\right\} = \cos(t)
\]
\end{resultado}

\subsection{Término 2: \texorpdfstring{$\Lap^{-1}\!\left\{\dfrac{1}{s(s^2+1)}\right\}$}{}}

Se reescribe la expresión como el producto $\dfrac{1}{s}\cdot\dfrac{1}{s^2+1}$ y se aplica el \textbf{teorema de convolución}. Identificando:
\[
    F(s) = \frac{1}{s} \;\Longrightarrow\; f(t) = 1,
    \qquad
    G(s) = \frac{1}{s^2+1} \;\Longrightarrow\; g(t) = \sin(t)
\]

Se calcula la convolución:
\[
    \Lap^{-1}\!\left\{\frac{1}{s(s^2+1)}\right\}
    = (1 * \sin)(t)
    = \int_0^t \sin(\tau)\,d\tau
    = \bigl[-\cos(\tau)\bigr]_0^t
    = 1 - \cos(t)
\]

\begin{resultado}
\[
    \Lap^{-1}\!\left\{\frac{1}{s(s^2+1)}\right\} = 1 - \cos(t)
\]
\end{resultado}

\subsection{Término 3: \texorpdfstring{$\Lap^{-1}\!\left\{\dfrac{s\,e^{-\frac{\pi}{2}s}}{(s^2+1)^2}\right\}$}{}}

Por el segundo teorema de traslación, se necesita primero calcular:
\[
    h(t) = \Lap^{-1}\!\left\{\frac{s}{(s^2+1)^2}\right\}
\]

\textbf{Aplicación del teorema de derivada en frecuencia.} Se observa que:
\[
    \frac{d}{ds}\left[\frac{1}{s^2+1}\right] = -\frac{2s}{(s^2+1)^2}
    \quad \Longrightarrow \quad
    \frac{s}{(s^2+1)^2} = -\frac{1}{2}\frac{d}{ds}\left[\frac{1}{s^2+1}\right]
\]

Usando el teorema de derivada con $n=1$ y $F(s) = \dfrac{1}{s^2+1}$, cuya inversa es $f(t) = \sin(t)$:
\[
    \Lap^{-1}\!\left\{-\frac{d}{ds}F(s)\right\} = t\,f(t) = t\sin(t)
\]

Por lo tanto:
\[
    \Lap^{-1}\!\left\{\frac{s}{(s^2+1)^2}\right\}
    = \frac{1}{2}\,t\sin(t)
\]

Aplicando ahora el segundo teorema de traslación con $a = \pi/2$ y usando que $\sin(t - \pi/2) = -\cos(t)$:
\begin{align*}
    \Lap^{-1}\!\left\{\frac{s\,e^{-\frac{\pi}{2}s}}{(s^2+1)^2}\right\}
    &= \frac{1}{2}\left(t-\frac{\pi}{2}\right)\sin\!\left(t-\frac{\pi}{2}\right)\,\U\!\left(t-\frac{\pi}{2}\right) \\[4pt]
    &= -\frac{1}{2}\left(t-\frac{\pi}{2}\right)\cos(t)\,\U\!\left(t-\frac{\pi}{2}\right)
\end{align*}

\begin{resultado}
\[
    \Lap^{-1}\!\left\{\frac{s\,e^{-\frac{\pi}{2}s}}{(s^2+1)^2}\right\}
    = -\frac{1}{2}\!\left(t-\frac{\pi}{2}\right)\cos(t)\,\U\!\left(t-\frac{\pi}{2}\right)
\]
\end{resultado}

\subsection{Término 4: \texorpdfstring{$\Lap^{-1}\!\left\{\dfrac{e^{-\frac{\pi}{2}s}}{s(s^2+1)}\right\}$}{}}

Por el segundo teorema de traslación y el resultado del Término~2:
\[
    \Lap^{-1}\!\left\{\frac{e^{-\frac{\pi}{2}s}}{s(s^2+1)}\right\}
    = \left[1 - \cos\!\left(t-\frac{\pi}{2}\right)\right]\U\!\left(t-\frac{\pi}{2}\right)
\]

Usando la identidad $\cos(t - \pi/2) = \sin(t)$:

\begin{resultado}
\[
    \Lap^{-1}\!\left\{\frac{e^{-\frac{\pi}{2}s}}{s(s^2+1)}\right\}
    = \bigl[1 - \sin(t)\bigr]\,\U\!\left(t-\frac{\pi}{2}\right)
\]
\end{resultado}

% ══════════════════════════════════════════════════════════════════════════════
\section{Solución final}
% ══════════════════════════════════════════════════════════════════════════════

Reuniendo todos los términos obtenidos:
\begin{align*}
    y(t) &= \cos(t)
           + \bigl(1 - \cos(t)\bigr)
           - \frac{1}{2}\!\left(t-\frac{\pi}{2}\right)\cos(t)\,\U\!\left(t-\frac{\pi}{2}\right)
           - \bigl[1 - \sin(t)\bigr]\,\U\!\left(t-\frac{\pi}{2}\right)
\end{align*}

Simplificando los términos sin escalón ($\cos t + 1 - \cos t = 1$) y factorizando el escalón:

\begin{resultado}
\[
    \boxed{
        y(t) = 1 + \U\!\left(t - \frac{\pi}{2}\right)
        \left[\sin(t) - 1 - \frac{1}{2}\!\left(t - \frac{\pi}{2}\right)\cos(t)\right]
    }
\]
\end{resultado}

\medskip

Esta solución puede verificarse evaluando en cada subintervalo:

\begin{itemize}
    \item Para $0 \leq t < \dfrac{\pi}{2}$: el escalón es nulo, por lo que $y(t) = 1$. \\
          Efectivamente, $y(0) = 1$ y $y'(t) = 0$ satisfacen las condiciones iniciales.

    \item Para $t \geq \dfrac{\pi}{2}$: el escalón se activa y la solución incorpora la respuesta al forzamiento $\sin(t)$, con la transición continua garantizada por la construcción del escalón en $t = \pi/2$.
\end{itemize}

% ══════════════════════════════════════════════════════════════════════════════
\end{document}